% Part: first-order-logic
% Chapter: introduction
% Section: formulas

\documentclass[../../../include/open-logic-section]{subfiles}

\begin{document}

\olfileid{fol}{int}{fml}

\section{\usetoken{P}{formula}}

We zullen !!{sentence}s van de eerste-ordelogica op de volgende wijze
rigoureus specificeren en de problemen behandelen die uit het gebruik
van !!{variable}s voortkomen. Eerst definiëren we een \emph{andere}
verzameling uitdrukkingen: !!{formula}s. Daarna kunnen we de rol van
!!{variable}s daarin onderzoeken. Met behulp van de begrippen
``vrije'' en ``gebonden'' !!{variable}s kunnen we vervolgens definiëren
wat !!a{sentence} is, namelijk !!a{formula} zonder vrije
!!{variable}s. Dit maakt niet alleen de definitie van ``!!{sentence}''
hanteerbaarder, maar is ook wezenlijk voor onze definitie van het
semantische begrip vervulling.

We definiëren ``!!{formula}'' eerst voor een eenvoudige taal van de
eerste orde. Deze bevat slechts één !!{predicate}~$\Obj P$, één
!!{constant}~$\Obj a$ en de logische symbolen $\lnot$, $\land$
en~$\lexists$. Onze volledige definities zullen veel algemener zijn:
we zullen oneindig veel !!{predicate}s en !!{constant}s toelaten. We
zullen bovendien !!{function}s beschouwen die met !!{constant}s en
!!{variable}s tot ``termen'' kunnen worden gecombineerd. Voorlopig zijn
$\Obj a$ en de variabelen onze enige termen. We hebben wel oneindig
veel !!{variable}s nodig. Officieel gebruiken we de symbolen
$\Obj v_0$, $\Obj v_1$, \dots, als variabelen.

\begin{defn}
De verzameling \emph{!!{formula}s}~$\Frm$ wordt als volgt gedefinieerd:
\begin{enumerate}
\item\ollabel{fmls-atom} $\Atom{\Obj P}{\Obj a}$ en $\Atom{\Obj
  P}{\Obj v_i}$ zijn !!{formula}s ($i \in \Nat$).

\tagitem{prvNot}{\ollabel{fmls-not}Als $!A$ !!a{formula} is, dan is $\lnot !A$
  !!{formula}.}{}

\tagitem{prvAnd}{Als $!A$ en $!B$ !!{formula}s zijn, dan is $(!A \land
  !B)$ !!a{formula}.}{}

\tagitem{prvEx}{\ollabel{fmls-ex}Als $!A$ !!a{formula} is en $x$ !!a{variable},
  dan is $\lexists[x][!A]$ !!a{formula}.}{}

\tagitem{limitClause}{\ollabel{fmls-limit}Niets anders is !!a{formula}.}{}
\end{enumerate}
\end{defn}

Uit \olref{fmls-atom} volgt dat $\Atom{\Obj P}{\Obj a}$ en
$\Atom{\Obj P}{\Obj v_i}$ voor elke $i \in \Nat$ !!{formula}s zijn.
Dit zijn de zogeheten \emph{atomaire} !!{formula}s. Zij vormen het
uitgangspunt. De overige clausules geven aan hoe uit reeds gevormde
!!{formula}s nieuwe !!{formula}s kunnen worden gevormd. Uit
\olref{fmls-not} volgt bijvoorbeeld dat $\lnot \Atom{\Obj P}{\Obj v_2}$
!!a{formula} is, omdat $\Atom{\Obj P}{\Obj v_2}$ volgens
\olref{fmls-atom} al !!a{formula} is. Vervolgens volgt uit
\olref{fmls-ex} dat $\lexists[\Obj v_2][\lnot \Atom{\Obj P}{\Obj
v_2}]$ eveneens !!{formula} is, enzovoort. \olref{fmls-limit} bepaalt
dat \emph{alleen} de op deze wijze gevormde tekenreeksen als
!!{formula}s gelden. In het bijzonder gelden $\lexists[\Obj
v_0][\Atom{\Obj P}{\Obj a}]$ en $\lexists[\Obj v_0][\lexists[\Obj
v_0][\Atom{\Obj P}{\Obj a}]]$ \emph{wel} als !!{formula}s, maar
$(\lnot \Atom{\Obj P}{\Obj a})$ niet, vanwege de overtollige buitenste
haakjes.

Deze wijze van definiëren van !!{formula}s heet een \emph{inductieve
definitie}. Zij maakt het mogelijk eigenschappen van !!{formula}s te
bewijzen met een vorm van inductiebewijs die \emph{structurele inductie}
heet. Beide begrippen worden algemeen besproken in
\olref[mth][ind][idf]{sec} en \olref[mth][ind][sti]{sec}; men dient die
paragrafen door te nemen voordat men zich in de latere bewijzen
verdiept. Het grondidee is als volgt. Wil men bewijzen dat iets voor
alle !!{formula}s geldt, dan toont men eerst aan dat het voor de
atomaire !!{formula}s geldt. Vervolgens toont men aan dat, \emph{als}
het voor willekeurige !!{formula}~$!A$ (en~$!B$) geldt, het \emph{ook}
geldt voor $\lnot !A$, $(!A \land !B)$ en $\lexists[x][!A]$. Hiermee
wordt bijvoorbeeld bewezen dat het voor $\lexists[\Obj v_2][\lnot
\Atom{\Obj P}{\Obj v_2}]$ geldt. Uit het eerste deel weten we dat het
voor de atomaire !!{formula}~$\Atom{\Obj P}{\Obj v_2}$ geldt. Uit het
tweede deel volgt dan dat het voor $\lnot \Atom{\Obj P}{\Obj v_2}$
geldt, en vervolgens ook voor $\lexists[\Obj v_2][\lnot \Atom{\Obj
P}{\Obj v_2}]$ zelf. Omdat alle !!{formula}s inductief uit atomaire
!!{formula}s worden voortgebracht, werkt dit voor elke formule.

\end{document}
